\section{Execution Validation and Illustrative Use}
\label{sec:validation}

We examine whether the implemented platform connects agent
decisions, environmental execution, and subsequent feedback as
specified. Execution checks and an individual trace assess
these connections, while aggregate outputs illustrate how the
platform supports studies of individual adjustment and
collective dynamics.


\subsection{Demonstration Setup}
\label{sec:demonstration-setup}

The demonstration uses an enriched Sioux Falls scenario with
16,822 simulated people and 11,793 fixed target trips. Full-day
activity chains and background mobility remain active in all
configurations. Table~\ref{tab:conditions} summarizes five
configurations that vary the target-trip decision mechanism
and the presence of an LLM manager. All use MATSim for
environmental execution.

\begin{table}[htbp]
    \centering
    \small
    \caption{Demonstration configurations. Trips outside the
    controlled target set retain MATSim behavioral mechanisms.
    MNL and LLM traveler agents select target actions from
    fixed route and departure-time sets.}
    \label{tab:conditions}
    \renewcommand{\arraystretch}{1.15}
    \begin{tabularx}{\linewidth}{
        @{}
        >{\raggedright\arraybackslash}p{0.29\linewidth}
        >{\raggedright\arraybackslash}X
        >{\raggedright\arraybackslash}p{0.20\linewidth}
        @{}
    }
        \toprule
        Configuration & Target-trip decisions & manager \\
        \midrule
        MATSim
        & MATSim plan selection and replanning
        & Disabled \\

        MNL
        & Constrained joint route--departure choice
        & Disabled \\

        LLM Traveler
        & LLM route--departure choice
        & Disabled \\

        MATSim + Manager
        & MATSim plan selection and replanning
        & Enabled \\

        LLM Traveler + Manager
        & LLM route--departure choice
        & Enabled \\
        \bottomrule
    \end{tabularx}
\end{table}

Each configuration completed 100 simulated days with simulation
seed 268. The comparison contains one run per configuration;
successive days share persistent states and are not independent
replications. Some executable builds differ, and the MNL
coefficients are uncalibrated. The results therefore describe
the recorded configurations rather than isolate individual
mechanisms. Detailed settings and measurement definitions are
provided in the appendices.

Day indices in this section follow the zero-based logs,
from Day0 to Day99. Pricing review opportunities occur on
Days20, 50, and 80. Positive charging begins on Day20 in
MATSim + Manager and on Day50 in LLM Traveler + Manager.
The figure labels Native, Traveler, Manager, and Joint
correspond to MATSim, LLM Traveler, MATSim + Manager,
and LLM Traveler + Manager, respectively.


\subsection{Execution and Feedback Consistency}
\label{sec:validation-execution}

All five configurations completed 11,793 target trips and 33,644 car legs per day, with no stuck agents. Work-arrival records cover the same 11,332 targets. Table~\ref{tab:execution-audit} separates checks of the combined run from the sampled inspection of decision inputs.

\begin{table}[tbp]
\centering\small
\caption{Execution audit scope and results. Full-record checks concern LLM Traveler + Manager; the information-access check covers the inspected samples only. Action checks apply to finalized outputs after validation and repair.}
\label{tab:execution-audit}
\renewcommand{\arraystretch}{1.15}
\begin{tabularx}{\linewidth}{@{}>{\raggedright\arraybackslash}p{0.25\linewidth} >{\raggedright\arraybackslash}p{0.29\linewidth} >{\raggedright\arraybackslash}X@{}}
\toprule
Check & Coverage & Recorded result \\
\midrule
Action constraints and execution & 1,179,300 finalized target decisions & Valid route and departure choices; accepted actions linked to applied plans and completed routes \\
Personal state continuity & 100-day personal memory chains & Continuous chains; finalized summaries within the configured limit \\
Pricing decisions and state & 3 policy decisions; 100 daily memory updates & Valid policy decisions; all daily memory updates completed \\
Information access & 4,000 sampled input contexts & No current announcements disclosed to uninformed travelers; no network-wide traffic information disclosed \\
\bottomrule
\end{tabularx}
\end{table}

Output repair affected 0.206\% of finalized travel-decision records and 0.250\% of personal-memory records. One manager-memory response required JSON repair. These results describe accepted outputs after validation and retry. Selected and realized departure times are recorded separately because preceding activities can delay execution. Appendix~\ref{app:audit-cost} reports exact repair counts and inference workload.

The audits check structured actions, recorded execution, and state continuity. Generated explanations and memory content require separate semantic checks, as the following cases illustrate.


\subsection{A Recorded Interaction Trace}
\label{sec:interaction-trace}

Figure~\ref{fig:execution-trace} follows synthetic traveler \texttt{10010\_2} in LLM Traveler + Manager. On Day~60, the agent selected route \texttt{c1} and the earliest available departure after receiving recent outcomes that included a late arrival on Day~59. The applied action and completed route matched the choice; the resulting outcome entered the rolling personal history for subsequent decisions.

\begin{figure}[!htbp]
\centering\small
\fbox{\begin{minipage}{0.94\linewidth}
\textbf{Observation before the Day~60 decision.} Four candidate routes, nine departure options, a 09:00 arrival reference, personal outcomes from Days~57--59, and current toll announcements. Network-wide conditions are unavailable. The Day~59 trip arrived at 09:10:48.\par\smallskip
\centerline{$\downarrow$}
\textbf{Structured action.} Route \texttt{c1}, departing at 05:56:43: the earliest available option and 90 minutes before the original plan.\par\smallskip
\centerline{$\downarrow$}
\textbf{Execution and recorded outcome.} Applied and realized actions agree with the record; the completed route matches \texttt{c1}. Car time is 21.97 minutes, destination arrival is 06:33:36, and payment is zero. Door-to-door time is 36.88 minutes and uses different boundaries from car time.\par\smallskip
\centerline{$\downarrow$}
\textbf{Subsequent decision.} On Day~61, the traveler selects \texttt{c0} at 05:56:43. The Day~60 outcome is available in the rolling personal history.
\end{minipage}}
\caption{A recorded interaction trace connecting permitted observations, a structured action, realized execution, and a subsequent decision. The case illustrates traceability; route, departure time, and shared conditions vary together.}
\label{fig:execution-trace}
\end{figure}

The case was selected to illustrate an inspectable decision sequence. Both route and departure time changed between Days~59 and 60, alongside changing environmental conditions, so the trace does not isolate their individual effects. A separate checked case illustrates an explanation inconsistency: traveler \texttt{10042\_2} on Day~61 selected a departure time of 07:46:19 and departed at that
time, while its generated explanation stated 07:49:39.. The executable action remained consistent with its record. These cases demonstrate what the archive can reveal; they do not estimate the prevalence of explanation errors.
\FloatBarrier

\subsection{Illustrative Collective Dynamics}
\label{sec:illustrative-use}

\paragraph{Individual and system-level outcomes.}

The configuration with the lowest population-wide car time
differs from the one with the shortest target-trip time. LLM Traveler records the lowest mean daily car time per
person, 19.90\% below MATSim, whereas LLM Traveler + Manager
records the shortest mean target-trip time. Relative to
LLM Traveler, the combined configuration has 6.63\% lower
target-trip time but 14.62\% higher mean daily car time
per person. Differences already exist before positive charging,
preventing attribution of this contrast to the manager
alone.

Figure~\ref{fig:outcomes} shows these trajectories across
different evaluation scopes. Mean daily car time per person
includes all car legs and waiting before network entry,
with total duration divided by the fixed population of
16,822 simulated people. Mean target-trip car time covers
the 11,793 controlled trips, while the work late-arrival
rate uses the fixed denominator of 11,332 work targets.
Together, these measures relate outcomes for controlled
trips to those across the full simulated population.

\begin{figure}[!htbp]
    \centering
    \includegraphics[
            width=\linewidth,
        trim= 0cm 0cm 0cm 0cm, clip
    ]{section/pic/outcomes.pdf}
    \caption{Individual and system-level mobility outcomes
across five configurations:
(a) mean daily car time per person (min/person);
(b) mean target-trip car time (min/trip); and
(c) work-target late-arrival rate (\%).
The denominators are 16,822 simulated people,
11,793 controlled trips, and 11,332 work targets,
respectively. Panel (a) includes all car legs and waiting
before network entry. Thin lines show daily values and
thick lines show trailing seven-day means.
Days~20 and 50 mark positive charging in Manager and Joint,
respectively; Day~90 marks the MATSim innovation cutoff.
Each configuration has one seed-268 run.}
    \label{fig:outcomes}
\end{figure}

\paragraph{Individual adjustment and aggregate patterns.}

The two LLM traveler configurations show relatively
persistent route choices alongside more frequent
departure-time adjustments. Across the 99 adjacent-day
transitions, LLM Traveler and the combined configuration change
routes in 3.08\% and 2.19\% of target decisions, respectively,
while changing departure times in 29.45\% and 28.04\%.
These overall rates include 2,780 targets with only one candidate
route. Among the 9,013 targets with multiple routes, the
corresponding route-switching rates are 4.03\% and 2.86\%.

\begin{figure}[htbp]
    \centering
    \includegraphics[
            width=\linewidth,
        trim= 0cm 0cm 0cm 0cm,
        clip
    ]{section/pic/behavior.pdf}
    \caption{Individual choice changes and aggregate route-flow
    turnover. (a) Population-wide route-switching rates across
    five configurations, shown on a logarithmic scale.
    (b) Departure-time adjustment rates for Traveler and Joint.
    (c) Individual route switching and net turnover of
    OD-specific route counts. Rates use all 11,793 targets
    over 99 adjacent-day transitions. Thin lines show daily
    values and thick lines show trailing seven-day means.
    Day50 marks Joint charging and Day90 marks the MATSim
    innovation cutoff.}
    \label{fig:behavior}
\end{figure}

Figure~\ref{fig:behavior} places individual changes alongside
aggregate route-flow turnover. Offsetting switches can leave
aggregate route counts relatively stable, making individual
adjustment and collective redistribution distinct measurements.
The decline in MATSim route switching near Day90 also reflects
the configured cutoff for route and activity-time innovation;
it cannot be interpreted solely as behavioral stabilization.

LLM Traveler and the combined configuration select departures
an average of 34.96 and 39.65 minutes earlier than the original
plans. Their time outcomes therefore need to be interpreted
together with schedule changes and the available action space.
These observations identify measurable dimensions of
agent-system dynamics without establishing learned habits,
coordination, or welfare improvements. Subsequent studies can
vary information, memory, and role interactions while retaining
the same execution and recording interfaces.

\FloatBarrier
